\documentclass[10pt,letterpaper]{article}

\usepackage[margin=0.72in]{geometry}
\usepackage[T1]{fontenc}
\usepackage{lmodern}
\usepackage{microtype}
\usepackage{parskip}
\usepackage{enumitem}
\usepackage{xcolor}
\usepackage{hyperref}

\definecolor{ink}{HTML}{18212B}
\definecolor{accent}{HTML}{0B5D68}
\definecolor{muted}{HTML}{52606D}
\hypersetup{colorlinks=true,urlcolor=accent,linkcolor=accent}
\pagestyle{empty}
\setlength{\parindent}{0pt}
\setlist[itemize]{leftmargin=1.35em,itemsep=0.18em,topsep=0.25em,parsep=0pt}
\setlength{\emergencystretch}{2em}

\newcommand{\sectiontitle}[1]{\vspace{0.65em}{\large\bfseries\color{accent}#1}\par\vspace{0.12em}\hrule\vspace{0.35em}}
\newcommand{\role}[2]{\textbf{#1}\hfill\textcolor{muted}{#2}\par}

\begin{document}
\color{ink}

\begin{center}
{\fontsize{24}{28}\selectfont\bfseries Jason Duran Dutton}\par
\vspace{0.12em}
{\large\color{accent}\bfseries Independent Researcher | Physics, Photonics, and Systems Engineering}\par
{\normalsize AMO / Optical Properties of Materials}\par
\vspace{0.35em}
Castlewood, Virginia, USA \textbar\ \href{mailto:sefidefi@outlook.com}{sefidefi@outlook.com} \textbar\ \href{mailto:jasondurandutton@pm.me}{jasondurandutton@pm.me}\par
\href{https://linkedin.com/in/jasondurandutton}{LinkedIn}\par
\href{https://github.com/JasonSEFIDEFI/SEFI-Manuscripts}{Research manuscripts} \textbar\ \href{https://github.com/JasonSEFIDEFI/SEFI_QEC_Stack}{Photonic QEC engine}
\end{center}

\sectiontitle{Profile}
Independent theoretical physics researcher with more than 15 years of self-directed study, manuscript development, and computational model building. My work connects geometric descriptions of electron and photon behavior with optical confinement, polarization, stability, and photonic error correction. I am particularly effective at turning difficult physics into precise, intuitive verbal explanations and at evaluating whether an explanation is physically coherent.

\textbf{Best fit for micro1:} expert-level spoken reasoning, physics explanation, conceptual error detection, and structured feedback on AI-generated answers in atomic, optical, and photonic physics.

\sectiontitle{AMO and Optical-Physics Focus}
\begin{itemize}
    \item \textbf{Atomic structure and transitions:} qualitative reasoning about hyperfine structure, polarization, selection-rule intuition, and the distinction between $\sigma$ and $\pi$ transitions.
    \item \textbf{Optical properties of materials:} refractive-index contrast, waveguide geometry, mode confinement, photonic stability, and the relationship between material response and field propagation.
    \item \textbf{Polarizability and external fields:} geometric interpretation of induced displacement, field response, and stability under perturbation.
    \item \textbf{Photonic systems:} SPDC, photon identity, CHSH/Bell behavior, displacement geometry, stabilizer logic, and geometric quantum error correction.
    \item \textbf{Model critique and explanation:} translating between physical intuition, mathematical structure, assumptions, and experimentally meaningful predictions.
\end{itemize}

\sectiontitle{Foundational and Theoretical Physics Focus}
\begin{itemize}
    \item \textbf{Geometric field modeling:} worldline geometry, curvature, torsion, invariants, field stability, and relationships between local structure and global behavior.
    \item \textbf{Quantum foundations:} electron/positron identity, interference, measurement, photon identity, SPDC, Bell/CHSH correlations, and interpretation of quantum state behavior.
    \item \textbf{Unified theoretical frameworks:} development of conceptual operators and layered models connecting spacetime geometry, identity-space structure, dynamical invariants, and physical interpretation.
    \item \textbf{Scientific reasoning:} separating assumptions, mathematical structure, physical intuition, predictions, and possible tests when analyzing a theory or technical explanation.
\end{itemize}

\sectiontitle{Computational, Systems, and Applied Research Focus}
\begin{itemize}
    \item \textbf{Scientific software and modeling:} Python-based research tools, modular simulation workflows, geometric state representation, error propagation, and reproducible technical documentation.
    \item \textbf{Quantum information systems:} geometric stabilizers, syndrome structure, displacement fields, coherence preservation, photonic encoding, and quantum error-correction concepts.
    \item \textbf{Industrial research translation:} converting research objectives into prints, controls, ladder logic, instrumentation, commissioning procedures, field tests, and operational decisions.
    \item \textbf{Reliability and safety engineering:} fault isolation, system integration, verification, regulatory awareness, high-consequence troubleshooting, and communication across technical disciplines.
\end{itemize}

\sectiontitle{Selected Research}
\role{Independent Theoretical Researcher}{2010--Present | Castlewood, VA}
\begin{itemize}
    \item Developed a unified geometric research program spanning worldline models, identity-space invariants, photonic stability, and quantum error correction.
    \item Built the \textbf{SEFI\_QEC\_Stack}, a photonic quantum error-correction engine for geometric stabilizers, syndrome structure, displacement fields, and correction mapping.
    \item Developed and documented interpretations of electron/positron identity, interference, measurement, SPDC, CHSH correlations, polarization evolution, and field stability.
    \item Prepared four research manuscripts for external review and made the supporting research portfolio publicly available for inspection.
    \item Refined physics explanations through iterative dialogue with AI systems, with emphasis on conceptual consistency, clear assumptions, and generalization across problems.
\end{itemize}

\sectiontitle{Research Portfolio}
\begin{itemize}
    \item \textbf{Geometric Worldline Foundations of the Electron Field} (JMP26-AR-02470): electron/positron identity, interference, and measurement.
    \item \textbf{SEFI: The Single Entity Field Interpretation} (JMP26-AR-02501): layered identity-space geometry and invariants.
    \item \textbf{Unified Geometric Field Theory from Worldline and Identity-Space Invariants} (JMP26-AR-02666): spacetime geometry and dynamical invariants.
    \item \textbf{Photonic Quantum Error Correction from Geometric Stability Principles} (ADV26-AR-04574): stability surfaces, displacement geometry, and correction mapping.
\end{itemize}

\sectiontitle{Technical and Communication Strengths}
\textbf{Physics:} AMO concepts \textbar\ optical response \textbar\ photonics \textbar\ quantum foundations \textbar\ geometric modeling\par
\textbf{Computing:} Python-based simulation workflows \textbar\ modular research software \textbar\ GitHub \textbar\ LaTeX\par
\textbf{Communication:} spoken technical explanation \textbar\ analogy with explicit limits \textbar\ assumption checking \textbar\ critique of model outputs

\sectiontitle{Professional Engineering Experience}
\role{Lead Electrician -- Pilot Process Research Installation}{Mineral Refining Company / Lone Mountain Refining | 2015--2016}
St. Charles, Virginia \textbar\ Virginia Tech-directed pilot deployment of patented Hydrophobic--Hydrophilic Separation (HHS) technology \textbar\ \href{https://mineralsrefining.com}{mineralsrefining.com}
\begin{itemize}
    \item Served as the first electrician on site and led full electrical integration for a pilot process designed to recover ultra-fine coal and metallic-ore particles from preparation-plant waste streams.
    \item Owned field coordination of control prints, ladder logic, sensor interfaces, system-level electrical decisions, installation verification, commissioning, and safety compliance.
    \item Integrated and tested phase-separation vessels, vibrating mixers, candle filters, pentane-based hydrophobic-liquid recovery loops, and advanced dewatering modules.
    \item Supported more than 1,200 hours of field testing under MSHA and DMME oversight, including validation of ultra-fine coal recovery below 45 micrometers and single-digit moisture control.
    \item Contributed to pilot results reporting feed-ash reduction from 40--67\% to 3--5\%, combustible recovery up to 97\%, and HHS applicability to metallic-ore particles below 10 micrometers and 1 micrometer.
    \item Supported environmental and economic outcomes including extended slurry-pond life, reduced waste deposition, improved plant economics, and recovery of previously impounded fines.
\end{itemize}

\role{Underground Electrical Repairman}{Paramount Contura / Alpha Metallurgical Resources | 2022--Present}
Safety-critical diagnostics and maintenance requiring disciplined troubleshooting, reliability-first decisions, and clear communication under operational constraints.

\role{Underground Electrician / Maintenance Foreman}{Alpha Natural Resources | 2010--2016}
Electrical systems maintenance, fault isolation, and safety compliance in a highly regulated industrial environment.

\role{Fabricator, Welder, and Foreman}{Dominion Steel | 2016--2018}
Led shop and field work while coordinating system interpretation, quality execution, and production deadlines.

\sectiontitle{Education and Qualification Context}
Castlewood High School -- High School Diploma, 2002\par
Mountain Empire Community College -- Underground and surface mining electrical study\par
Brevard Community College / UCF -- CIS and business-management coursework\par
My physics preparation is an independent, sustained research program rather than an awarded PhD; the linked manuscripts and software provide the most direct evidence of my current research level and technical scope.

\vfill
\begin{center}
\small\color{muted}Portfolio: \href{https://github.com/JasonSEFIDEFI/SEFI-Manuscripts}{github.com/JasonSEFIDEFI/SEFI-Manuscripts}
\end{center}

\end{document}
